%==============================================================================
% guion_defensa.tex — Guion (script) para la defensa oral del trabajo de grado
% Tratamiento Numérico de la KHI en RRMHD
%
% Bryan Martinez Anzola & Sebastián Rodriguez García
% Director: Sergio Miranda-Aranguren — UD Francisco José de Caldas
% Mayo 2026 — Duración objetivo: 45 minutos
%
% COMPILACIÓN:  pdflatex guion_defensa.tex   (no requiere bibliografía)
%
% NOTAS DE USO:
%  * Cada bloque corresponde a una diapositiva. La cabecera indica el orador
%    asignado, número de diapositiva y duración estimada (a ritmo conservador
%    de ~150 wpm; a ritmo de defensa típico ~170-180 wpm el tiempo total se
%    aproxima a los 45 min objetivo).
%  * Las marcas «[ENLACE → ...]» son señales orales explícitas para hacer el
%    relevo entre presentadores con naturalidad.
%  * El estilo es formal-académico y técnicamente preciso; las cifras y
%    parámetros son consistentes con la presentación entregada.
%==============================================================================
\documentclass[11pt,a4paper]{article}

\usepackage[utf8]{inputenc}
\usepackage[T1]{fontenc}
% babel: provide=* permite usar el idioma aunque la distribución mínima
% no traiga el .ldf clásico de español. Esta línea es robusta tanto en
% TeX Live completo (texlive-lang-spanish) como en instalaciones mínimas.
\usepackage[spanish,provide=*]{babel}
\IfFileExists{lmodern.sty}{\usepackage{lmodern}}{}
\IfFileExists{microtype.sty}{\usepackage[expansion=false]{microtype}}{}
\usepackage[margin=2.3cm]{geometry}
\usepackage{xcolor}
\usepackage{tcolorbox}
\tcbuselibrary{skins,breakable}
\usepackage{enumitem}
\usepackage{fancyhdr}
\usepackage{titlesec}
\usepackage{amsmath,amssymb}
\usepackage{hyperref}

% ---- Paleta de colores diferenciada por orador --------------------------------
\definecolor{seb}{HTML}{1F4E79}   % Azul corporativo — Sebastián
\definecolor{bry}{HTML}{8A2A12}   % Rojo terracota — Bryan
\definecolor{neutral}{HTML}{2E2E2E}
\definecolor{tagbg}{HTML}{F2F2F2}

\hypersetup{colorlinks=true,linkcolor=neutral,urlcolor=seb}

% ---- Cabeceras y pie de página ------------------------------------------------
\pagestyle{fancy}
\fancyhf{}
\fancyhead[L]{\small Defensa --- KHI en RRMHD}
\fancyhead[R]{\small Mayo 2026}
\fancyfoot[C]{\small \thepage}
\renewcommand{\headrulewidth}{0.4pt}

% ---- Estilo de títulos --------------------------------------------------------
\titleformat{\section}
  {\Large\bfseries\color{neutral}}{\thesection}{0.7em}{}
\titlespacing*{\section}{0pt}{1.3em}{0.5em}

% ---- Bloques de orador --------------------------------------------------------
% Sintaxis: \turnSeb{slide}{título}{tiempo}{acumulado}{ texto }
%           \turnBry{slide}{título}{tiempo}{acumulado}{ texto }
\newtcolorbox{seborador}[4]{%
  enhanced, breakable,
  colback=white, colframe=seb!75!black,
  boxrule=0.6pt, arc=2pt,
  left=8pt, right=8pt, top=6pt, bottom=6pt,
  title={\bfseries\color{white}
         \textsc{Sebastián} \ \textbar\ \ Diap.~#1 \ \textbar\ \ #2
         \hfill \small #3 \ (acum.\ #4)},
  coltitle=white, colbacktitle=seb,
  fonttitle=\small\bfseries
}
\newtcolorbox{bryorador}[4]{%
  enhanced, breakable,
  colback=white, colframe=bry!85!black,
  boxrule=0.6pt, arc=2pt,
  left=8pt, right=8pt, top=6pt, bottom=6pt,
  title={\bfseries\color{white}
         \textsc{Bryan} \ \textbar\ \ Diap.~#1 \ \textbar\ \ #2
         \hfill \small #3 \ (acum.\ #4)},
  coltitle=white, colbacktitle=bry,
  fonttitle=\small\bfseries
}

% Macro para resaltar el momento de relevo entre oradores
\newcommand{\relevo}[1]{%
  \par\medskip
  \noindent\colorbox{tagbg}{\parbox{\dimexpr\linewidth-2\fboxsep\relax}{%
    \small\itshape \textbf{[Relevo]} #1}}\par\medskip}

% Macro para nota escénica (pausa, entonación)
\newcommand{\escena}[1]{\textit{\textcolor{neutral!70!white}{(#1)}}}

\setlength{\parskip}{0.4em}
\setlength{\parindent}{0pt}

%==============================================================================
\begin{document}

% ---- Portada del guion --------------------------------------------------------
\begin{titlepage}
\centering
\vspace*{2cm}
{\Large\bfseries\color{neutral} Guion para la Defensa Oral}\\[0.4em]
{\large Trabajo de Grado --- Programa Académico de Física}\\[1.5em]
{\Huge\bfseries Tratamiento Numérico de la}\\[0.3em]
{\Huge\bfseries Inestabilidad Kelvin--Helmholtz}\\[0.3em]
{\Large bajo el marco de la}\\[0.3em]
{\Huge\bfseries Magnetohidrodinámica Relativista Resistiva}\\[2em]

\rule{0.6\linewidth}{0.4pt}\\[1em]
{\large
  \textbf{Bryan Martinez Anzola} \quad $\cdot$ \quad
  \textbf{Sebastián Rodriguez García}\\[0.5em]
  Director: Prof.\ Ph.D.\ Sergio Miranda-Aranguren\\[0.3em]
  Universidad Distrital Francisco José de Caldas\\[0.3em]
  Bogotá, Mayo de 2026
}\\[1.5em]
\rule{0.6\linewidth}{0.4pt}\\[2em]

\begin{tcolorbox}[colback=tagbg, colframe=neutral!50, width=0.85\linewidth,
                  arc=2pt, boxrule=0.5pt]
\small
\textbf{Objetivo de duración:} 45 minutos.\\
\textbf{Distribución de oradores:}\\[0.3em]
\textcolor{seb}{\rule{0.8em}{0.8em}}\ \textsc{Sebastián}: \ 1, 2, 7, 8, 10, 11, 12, 15, 16, 20, 22, 23, 24, 28, 29, 33, 34, 35\\[0.2em]
\textcolor{bry}{\rule{0.8em}{0.8em}}\ \textsc{Bryan}: \ 3, 4, 5, 6, 9, 13, 14, 17, 18, 19, 21, 25, 26, 27, 30, 31, 32, 36, 37
\end{tcolorbox}
\vfill
\end{titlepage}

\tableofcontents
\newpage

%==============================================================================
\section{Apertura --- Diapositivas 1--2 (Sebastián)}
%==============================================================================

\begin{seborador}{1}{Portada}{0:30}{0:30}
Buenas tardes, profesores y compañeros. Mi nombre es Sebastián Rodríguez
García y, junto con mi compañero Bryan Martínez Anzola y bajo la dirección
del profesor Sergio Miranda-Aranguren, presentamos hoy el avance de nuestro
trabajo de grado titulado «Tratamiento numérico de la inestabilidad
Kelvin--Helmholtz en el marco de la magnetohidrodinámica relativista
resistiva», desarrollado sobre el código \textsc{cueva}. A lo largo de los
próximos cuarenta y cinco minutos expondremos el origen documental del
estudio, el marco teórico que lo sustenta, la metodología numérica empleada
y los resultados cuantitativos de nuestra campaña de cuarenta y cuatro
simulaciones validadas. \escena{breve pausa, mirada al jurado}
\end{seborador}

\begin{seborador}{2}{Estructura de la presentación}{0:50}{1:20}
La presentación está organizada en nueve bloques. Iniciaremos con el origen
documental del estudio, que es un póster elaborado el año anterior en el
laboratorio de astrofísica computacional de la universidad. Luego revisaremos
la propuesta formal de tesis con su justificación astrofísica. El núcleo
teórico se desarrolla en tres bloques: electrodinámica covariante, tensor
energía--momento con el sistema GLM, y ley de Ohm relativista con el sistema
conservativo. Posteriormente abordaremos la inestabilidad Kelvin--Helmholtz
desde la formulación clásica, pasando por la magnetohidrodinámica ideal y
relativista, hasta su descripción en RRMHD. Cerraremos con el avance en la
monografía, la campaña de simulaciones y el estado actual del proyecto.
\relevo{Bryan continuará con el bloque que origina toda la investigación.}
\end{seborador}

%==============================================================================
\section{Origen del estudio --- Diapositivas 3--6 (Bryan)}
%==============================================================================

\begin{bryorador}{3}{Origen Documental: Póster del Laboratorio}{1:15}{2:35}
Gracias, Sebastián. Todo este trabajo nace del póster que ven en pantalla,
presentado el año pasado en el laboratorio de astrofísica computacional
dirigido por el profesor Miranda-Aranguren. La mitad superior contiene la
formulación inicial de la inestabilidad Kelvin--Helmholtz bajo el sistema
RRMHD, las condiciones iniciales del campo de velocidad con el perfil
tangente hiperbólico y la configuración del campo magnético con polaridades
opuestas a cada lado de la interfaz. La mitad inferior muestra los resultados
preliminares: campos de presión, densidad y componentes de velocidad para
distintos valores de la conductividad eléctrica $\sigma$. En esa exploración
inicial ya se intuía cualitativamente que la resistividad afecta el desarrollo
de la inestabilidad, pero no había una caracterización cuantitativa del
fenómeno. Este póster constituye el punto de partida documental del trabajo,
y precisamente sobre las limitaciones que identificamos allí se construyó la
propuesta formal de tesis.
\end{bryorador}

\begin{bryorador}{4}{Punto de Partida: Lecciones del Póster}{1:30}{4:05}
Al revisar el póster identificamos cuatro limitaciones críticas. Primero, las
simulaciones preliminares no alcanzaban el régimen no lineal de la
inestabilidad: se quedaban atrapadas en la fase de crecimiento exponencial sin
llegar al enrollamiento ni a la cascada turbulenta. Segundo, la resolución
espacial era insuficiente para resolver la capa de cizalladura con el detalle
necesario. Tercero, el costo computacional resultaba prohibitivo para hacer
un barrido sistemático en $\sigma$. Y cuarto, no existían diagnósticos
cuantitativos: solo había visualizaciones cualitativas de la vorticidad. La
configuración usaba el perfil tangente hiperbólico, campos magnéticos de
polaridad opuesta y un rango exploratorio de $\sigma$ entre cien y diez mil.
El código \textsc{cueva}, desarrollado por nuestro director, ya contaba con
el solucionador HLLC, los esquemas IMEX Runge--Kutta y la limpieza de
divergencias GLM. La lección clave fue clara: la barrera hacia el régimen no
lineal era a la vez \emph{computacional} y \emph{metodológica}. Mientras el
póster se quedó en la visualización cualitativa de $\omega_z$, la tesis se
propuso entregar $\gamma_\text{KHI}$ como función de $\sigma$, la enstrofía
perturbada $\Omega'_z(t)$, $\sigma_\text{crit}$, una ley de potencias y un
modelo paramétrico de la transición.
\end{bryorador}

\begin{bryorador}{5}{Régimen Lineal vs.\ No Lineal}{1:25}{5:30}
Esta diapositiva formaliza la frontera que el póster no pudo cruzar. En el
régimen lineal las perturbaciones son pequeñas frente al estado base, lo que
permite descomponer las variables como $\mathbf{U}=\mathbf{U}_0+\delta
\mathbf{U}\,e^{i(\mathbf{k}\cdot\mathbf{x}-\omega t)}$. La amplitud crece como
$e^{\gamma_\text{KHI}t}$. Este régimen sí lo alcanzamos en el póster y es el
que tiene predicciones analíticas en la literatura, útiles para validar el
código. Sin embargo, no captura los fenómenos de mayor interés astrofísico:
reconexión, formación de vórtices y mezcla turbulenta. La meta de la tesis es
justamente ese régimen no lineal: saturación de la fase exponencial, formación
de vórtices de Kelvin--Helmholtz, reconexión magnética, conversión de energía
cinética a magnética y desorden topológico de las líneas de campo. El
observable que utilizamos es la enstrofía de la perturbación: como $\Omega'_z$
es proporcional al cuadrado de la amplitud, su logaritmo crece linealmente en
el tiempo con pendiente $2\gamma_\text{KHI}$. La gráfica a la derecha muestra
$\ln\Omega'_z(t)$ para las cuarenta y cuatro simulaciones validadas, con la
ventana lineal $t\in[2.4,3.4]$ destacada en amarillo y un coeficiente de
determinación promedio de $\bar R^2=0.92$.
\end{bryorador}

\begin{bryorador}{6}{La Propuesta: Título, Objetivo, Hipótesis}{1:15}{6:45}
Esta diapositiva consolida la propuesta formal entregada en diciembre del
año pasado. El título es el ya enunciado y el objetivo general es simular
numéricamente el desarrollo de la inestabilidad Kelvin--Helmholtz en plasmas
relativistas usando el marco RRMHD con el código \textsc{cueva}. Los objetivos
específicos son cuatro: cuantificar cómo $\sigma$ modifica
$\gamma_\text{KHI}$; analizar la influencia de la resistividad en la formación
de estructuras secundarias; identificar el efecto de la difusión magnética
sobre las variables globales del sistema; y evaluar el rol de la orientación e
intensidad del campo magnético en la estabilización. La hipótesis central es
que la resistividad física actúa como mecanismo estabilizador: al aumentar
$\eta=1/\sigma$ se reduce la tensión magnética efectiva e incrementa la
difusión, disminuyendo $\gamma_\text{KHI}$. El cronograma se dividió en dos
semestres: el primero dedicado a formulación y validación, el segundo a
simulaciones extensivas y a la escritura del documento.
\relevo{Le devuelvo la palabra a Sebastián para la justificación astrofísica.}
\end{bryorador}

%==============================================================================
\section{Justificación astrofísica --- Diapositivas 7--8 (Sebastián)}
%==============================================================================

\begin{seborador}{7}{Justificación Astrofísica I}{1:30}{8:15}
Gracias, Bryan. La justificación de este trabajo se sostiene sobre una
observación fundamental: más del noventa por ciento de la materia bariónica
del universo se encuentra en estado de plasma. La magnetohidrodinámica es,
por tanto, la herramienta natural para describir entornos de alta energía
como núcleos galácticos activos, chorros relativistas, vientos de púlsares,
brotes de rayos gamma y magnetares. La inestabilidad Kelvin--Helmholtz
aparece de manera ubicua en todos estos sistemas porque, siempre que existe
un gradiente fuerte de velocidad en un fluido magnetizado, la inestabilidad
emerge. Y emerge para hacer cosas físicamente relevantes: genera mezcla entre
capas, amplifica el campo magnético por estiramiento de líneas, induce
reconexión magnética turbulenta y redistribuye masa, momento y energía.
La pregunta inmediata es por qué necesitamos resistividad. La RMHD ideal
asume $\sigma\to\infty$, lo que congela las líneas de campo al fluido y
prohíbe físicamente la reconexión. Cuando un código RMHD ideal igualmente
muestra reconexión, ésta es de origen \emph{numérico no controlado},
dependiente del esquema. La RRMHD con $\sigma$ finita permite reconexión
predecible y física, disipación realista por efecto Joule, modelado de capas
de corriente y control explícito de la difusión magnética.
\end{seborador}

\begin{seborador}{8}{Justificación Astrofísica II}{1:15}{9:30}
Profundizando en la transición, la limitación del RMHD ideal es severa:
invalida cuantitativamente los resultados en llamaradas de núcleos activos,
en magnetosferas de púlsares y en cualquier región con capas de corriente,
que es justamente donde la KHI es más interesante. Al incluir el término
resistivo $\eta\nabla^2\mathbf{B}$, dominante en regiones de altos
gradientes, recuperamos aniquilación magnética controlada, conversión de
energía magnética en calor y efectos dinámicos sobre las partículas. El
contraste se ve nítido en la ley de Ohm: en MHD ideal el campo eléctrico
comóvil es estrictamente cero y la topología de $\mathbf{B}$ es invariante;
en RRMHD la corriente tiene una contribución resistiva
$\sigma\Gamma[\mathbf{E}+\mathbf{v}\times\mathbf{B}-(\mathbf{E}\cdot
\mathbf{v})\mathbf{v}]$ que permite difusión controlada y reconexión física.
Las referencias clave del proyecto son Pimentel y Lora-Clavijo de 2019,
quienes demostraron que la resistividad explícita altera fundamentalmente
la forma y la evolución no lineal de la KHI; y Mizuno de 2013, que mostró
cómo la ecuación de estado modula el calentamiento resistivo del plasma.
\relevo{Le paso la palabra a Bryan para abrir el marco teórico.}
\end{seborador}

%==============================================================================
\section{Marco teórico --- Diapositiva 9 (Bryan)}
%==============================================================================

\begin{bryorador}{9}{El Sistema RRMHD: Visión General}{1:30}{11:00}
Gracias, Sebastián. El marco RRMHD se construye sobre cinco ingredientes
fundamentales que enumero brevemente: la conservación de masa, expresada como
la divergencia covariante de $\rho u^\mu$ igual a cero; la conservación de
energía--momento, divergencia del tensor $T^{\mu\nu}$ igual a cero; las
ecuaciones de Maxwell en sus pares homogéneo e inhomogéneo; la ley de Ohm
relativista, que cierra el sistema acoplando la corriente con el campo
electromagnético en el marco comóvil; y, finalmente, el sistema GLM aumentado,
que controla las divergencias numéricas espurias.

Trabajamos en unidades naturales racionalizadas Heaviside--Lorentz, con
$c=\mu_0=\epsilon_0=1$, métrica de Minkowski con signatura
$(-,+,+,+)$ y la convención usual de índices griegos para cuatro dimensiones
e índices latinos para tres. El vector de estado conservativo de
\textsc{cueva} tiene trece componentes independientes: la densidad bariónica
$D$, las tres componentes del momento $S^j$, la energía $\tau$, las tres
componentes de $\mathbf{B}$, las tres de $\mathbf{E}$ y los dos campos
auxiliares $\Psi$ y $\Phi$ del esquema GLM. Las variables cinemáticas se
escriben con el factor de Lorentz $\Gamma=(1-v^2)^{-1/2}$.
\relevo{Le paso la palabra a Sebastián para entrar en el detalle del electromagnetismo covariante.}
\end{bryorador}

%==============================================================================
\section{Electrodinámica covariante --- Diapositivas 10--12 (Sebastián)}
%==============================================================================

\begin{seborador}{10}{Cuadripotencial y Tensor de Faraday}{1:50}{12:50}
Tomo la palabra para desarrollar la electrodinámica desde su formulación
geométrica más limpia: las formas diferenciales. El cuadripotencial $A$ se
trata como una uno-forma diferencial sobre el fibrado cotangente del
espacio-tiempo, con componentes $A_\mu=(-\phi,\mathbf{A})$. La invariancia
gauge significa que dos potenciales que difieren en $d\Lambda$, derivada
exterior de una función escalar, producen exactamente los mismos observables.
La derivada exterior es nilpotente: $d^2=0$. Esto implica directamente que el
tensor de Faraday $F=dA$ es estrictamente invariante de gauge, como puede
leerse en el panel superior derecho. Cuando hago la derivada exterior
explícitamente, sustituyendo $A=A_\nu\,dx^\nu$ y aplicando Leibniz junto con
la antisimetría del producto cuña, obtengo la expresión usual del tensor de
Faraday: $F_{\mu\nu}=\partial_\mu A_\nu-\partial_\nu A_\mu$. Sus componentes
en notación matricial son las que aparecen abajo: la fila cero contiene las
componentes del campo eléctrico cambiadas de signo, y el bloque tres por tres
espacial recoge las componentes del campo magnético en arreglo antisimétrico.
Igualando esta forma matricial con la expresión variacional, recuperamos las
relaciones clásicas: $\mathbf{E}=-\nabla\phi-\partial_t\mathbf{A}$ y
$\mathbf{B}=\nabla\times\mathbf{A}$. La dos-forma $F$ es entonces el andamiaje
exacto del electromagnetismo clásico, derivado por pura geometría diferencial.
\end{seborador}

\begin{seborador}{11}{Descomposición Relativa a un Observador}{1:35}{14:25}
En esta diapositiva proyectamos el tensor de Faraday sobre el marco comóvil
del fluido. Tomamos un observador con cuadrivelocidad $u^\mu$, normalizada
como $u^\mu u_\mu=-1$. El campo eléctrico comóvil se obtiene contrayendo
$F^{\mu\nu}$ con $u_\nu$ y resulta ortogonal a $u^\mu$, lo que confirma que es
puramente espacial en el marco propio. El campo magnético comóvil se obtiene
de la misma manera pero usando el tensor dual ${}^*F^{\mu\nu}$, definido con
el tensor totalmente antisimétrico de Levi-Civita. Ambos vectores tienen tres
grados de libertad cada uno, completando los seis grados de libertad del
campo electromagnético. Con esta descomposición, el tensor de Faraday se
reconstruye como aparece en el panel central.

La importancia física es enorme: en el límite MHD ideal, $\sigma\to\infty$, y
para evitar corrientes divergentes el campo eléctrico comóvil debe ser cero.
Esto colapsa el tensor a una componente puramente magnética, lo que es la
formulación rigurosa del congelamiento topológico. En el régimen RRMHD, en
cambio, $E^\mu\neq 0$, lo que permite calentamiento de Joule
($j^\mu E_\mu>0$), difusión del campo magnético y cambios reales de
topología. Cierro mencionando el cuadrivector de Poynting, que captura el
flujo de energía electromagnética en forma covariante.
\end{seborador}

\begin{seborador}{12}{Ecuaciones de Maxwell en Forma Covariante}{1:25}{15:50}
En esta última diapositiva del bloque que me corresponde, condenso las
ecuaciones de Maxwell en sus dos pares. El par homogéneo emerge de la
nilpotencia: como $F=dA$, su derivada exterior es automáticamente cero. Esto
contiene simultáneamente la ausencia de monopolos magnéticos y la ley de
Faraday. El origen es puramente geométrico, no requiere fuentes y prohíbe los
monopolos magnéticos aislados. El par inhomogéneo se obtiene del principio de
mínima acción aplicado al lagrangiano electromagnético acoplado a una
corriente: $\partial_\mu F^{\mu\nu}=J^\nu$.

Lo notable: la conservación de la carga no se postula, emerge automáticamente
de la antisimetría del tensor. Aplicando dos derivadas a $F^{\mu\nu}$
obtenemos cero por construcción, lo que implica $\partial_\nu J^\nu=0$. El
significado geométrico de los productos cuña $dx^\mu\wedge dx^\nu$ es que
representan elementos de área orientados en el espacio-tiempo: $dx\wedge dy$
son superficies de flujo magnético, $dt\wedge dx$ son áreas de circulación
eléctrica. La base diferencial absorbe contracciones de Lorentz y eso
garantiza la covariancia relativista.
\relevo{Le devuelvo la palabra a Bryan para el tensor de energía--momento.}
\end{seborador}

%==============================================================================
\section{Tensor de energía--momento --- Diapositivas 13--14 (Bryan)}
%==============================================================================

\begin{bryorador}{13}{Tensor Energía-Momento: Derivación Variacional I}{1:50}{17:40}
Gracias, Sebastián. El tensor de energía--momento se construye desde el
principio variacional, como propone Hilbert. Para una acción dada,
$T^{\mu\nu}$ se define como $-\frac{2}{\sqrt{-g}}\,\frac{\delta S}
{\delta g_{\mu\nu}}$. En espacio plano de Minkowski, $g_{\mu\nu}\to
\eta_{\mu\nu}$ y $\sqrt{-g}\to 1$. La aditividad es esencial: el tensor total
es la suma del tensor del fluido perfecto y el tensor electromagnético.

Para el campo electromagnético, el lagrangiano es
$\mathcal{L}_\text{em}=-\tfrac{1}{4}F_{\alpha\beta}F^{\alpha\beta}$. Aplicando
el operador de Hilbert obtenemos el tensor de Maxwell:
$T^{\mu\nu}_\text{em}=F^{\mu\lambda}F^\nu{}_\lambda-\tfrac{1}{4}\eta^{\mu\nu}
F_{\alpha\beta}F^{\alpha\beta}$. Es simétrico y tiene traza nula en vacío, lo
cual es huella de la invariancia conforme del electromagnetismo libre. Para
el fluido perfecto, $\mathcal{L}_\text{fluid}=p$, y se obtiene el tensor
habitual: $\rho h\,u^\mu u^\nu+p\,\eta^{\mu\nu}$, con entalpía específica
$\rho h=\epsilon+p$.

Aún más interesante: el tensor electromagnético admite una descomposición
covariante muy compacta en términos de $E^\mu$, $B^\mu$ y el cuadrivector de
Poynting $S^\mu$, donde $h^{\mu\nu}=\eta^{\mu\nu}+u^\mu u^\nu$ es el
proyector espacial. Esta forma es la que aparece directamente implementada
en el código \textsc{cueva} y la que permite identificar los flujos físicos
cuando hagamos la descomposición $3+1$.
\end{bryorador}

\begin{bryorador}{14}{Tensor Energía-Momento: Derivación Variacional II}{1:40}{19:20}
En esta segunda parte de la derivación variacional llegamos a la dinámica
acoplada materia--campo. La conservación local del tensor total exige que
cualquier variación en la energía o momento del campo electromagnético se
compense exactamente con una variación opuesta en el fluido. Calculamos la
divergencia del tensor electromagnético usando la ley inhomogénea de Maxwell
y la identidad de Bianchi, y obtenemos un resultado clave: la divergencia de
$T^{\mu\nu}_\text{em}$ es $-F^{\nu\lambda}J_\lambda$. Sustituyendo en la
conservación total, la divergencia del tensor del fluido es
$F^{\nu\lambda}J_\lambda$, es decir, la generalización relativista directa
de Euler--Navier--Stokes con forzamiento electromagnético. Lo notable es que
\emph{la fuerza de Lorentz no se postula}: emerge automáticamente de la
geometría del tensor de Faraday y de la conservación covariante.

Las densidades conservadas observables en el laboratorio son: la energía
total, suma del término electromagnético $\tfrac{1}{2}(E^2+B^2)$ y el
hidrodinámico $\rho h\Gamma^2-p$; el momento total, suma del vector de
Poynting $\mathbf{E}\times\mathbf{B}$ y el momento hidrodinámico
$\rho h\Gamma^2 \mathbf{v}$; y la masa bariónica $D=\rho\Gamma$. Finalmente,
este sistema covariante no es directamente integrable numéricamente:
necesitamos distinguir evolución temporal de gradientes espaciales mediante
la descomposición $3+1$, proyectando sobre las hipersuperficies espaciales
del observador euleriano.
\relevo{Le paso la palabra a Sebastián para el sistema GLM.}
\end{bryorador}

%==============================================================================
\section{Sistema GLM --- Diapositivas 15--16 (Sebastián)}
%==============================================================================

\begin{seborador}{15}{Sistema GLM Aumentado I}{1:40}{21:00}
Tomo la palabra para abordar uno de los puntos más sutiles de toda la
formulación: el sistema GLM. Las restricciones $\nabla\cdot\mathbf{B}=0$ y
$\nabla\cdot\mathbf{E}=\rho_e$ son ecuaciones \emph{elípticas}, sin derivada
temporal. Esto las hace numéricamente patológicas: los errores de
truncamiento del esquema generan monopolos magnéticos espurios que no se
propagan ni se disipan, sino que se acumulan estáticamente y terminan
destruyendo la estabilidad del cálculo, especialmente en simulaciones largas
como las nuestras.

La solución, propuesta originalmente por Dedner y colaboradores en 2002, es
introducir multiplicadores de Lagrange generalizados, dos campos escalares
$\Phi$ y $\Psi$, que convierten las restricciones elípticas en ecuaciones
dinámicas. Los errores se propagan a velocidad finita y se disipan, en lugar
de acumularse, sin alterar las soluciones físicas exactas. Esto puede
formalizarse como una acción extendida: incluyendo un término de tipo
Stueckelberg $-\Phi\,\partial_\mu A^\mu$, recuperamos por variación dos
cosas: la ecuación de movimiento modificada con un término gradiente de
$\Phi$ en la corriente, y una ecuación de onda para $\Phi$ cuya fuente es la
divergencia de $A$. La justificación física es elegante: desde la TCC, el
campo $\Phi$ opera como un mecanismo de Stueckelberg, aislando los modos
longitudinales erróneos generados por la discretización e impidiendo que
contaminen los modos transversales físicos.
\end{seborador}

\begin{seborador}{16}{Sistema GLM Aumentado II}{1:40}{22:40}
En la segunda parte del bloque GLM completamos el sistema con el campo dual
$\Psi$, asociado a la ley de Bianchi. Por la simetría de Hodge de las
ecuaciones de Maxwell, el dual del tensor de Faraday admite el mismo
tratamiento, lo que da la ecuación de campo $\Psi$ mostrada arriba a la
izquierda. El sistema GLM covariante completo incluye entonces cuatro
ecuaciones: las dos modificadas de Maxwell y las dos ecuaciones de evolución
para los campos auxiliares con sus respectivos coeficientes de
amortiguamiento $\kappa_E$ y $\kappa_B$. La descomposición $3+1$ produce las
cuatro ecuaciones eulerianas que aparecen abajo a la izquierda y que son las
efectivamente integradas en \textsc{cueva}.

El elemento más bonito aquí es la ecuación del telégrafo: los errores
espurios $\xi$ obedecen $\partial_t^2\xi+\kappa\partial_t\xi-\nabla^2\xi=0$,
cuya solución envolvente decae exponencialmente como $e^{-\kappa t/2}$. Esto
da el comportamiento dual que justifica el método: en el límite $\kappa\to 0$
recuperamos propagación hiperbólica pura a velocidad uno, conservativa y
causal pero no disipativa. En el límite $\kappa\to\infty$ recuperamos una
ecuación de difusión, eficaz pero con condición de Courant prohibitiva
$\Delta t\propto\Delta x^2$. El régimen GLM es justamente el intermedio:
mantiene causalidad relativista, mantiene el paso de tiempo proporcional a
$\Delta x$, y disipa los errores con decaimiento exponencial.
\relevo{Le paso la palabra a Bryan para la ley de Ohm.}
\end{seborador}

%==============================================================================
\section{Ley de Ohm y sistema conservativo --- Diapositivas 17--19 (Bryan)}
%==============================================================================

\begin{bryorador}{17}{La Ley de Ohm Relativista I}{1:35}{24:15}
Gracias, Sebastián. La ley de Ohm es el ingrediente que cierra el sistema y
diferencia el RRMHD del RMHD ideal. Empezamos por la descomposición ortogonal
de la cuadricorriente: usando el proyector espacial
$h^\mu{}_\nu=\delta^\mu_\nu+u^\mu u_\nu$, separamos la corriente en una parte
convectiva $\rho_e u^\mu$, paralela a la cuadrivelocidad, y una parte
conductiva $j^\mu$, ortogonal a ella. La densidad de carga propia
$\rho_e\equiv -J^\nu u_\nu$ es un escalar invariante de Lorentz. La corriente
de conducción $j^\mu\equiv h^\mu{}_\nu J^\nu$ es también invariante.

En el régimen lineal e isotrópico ---el nuestro, ignorando efectos Hall y
termoeléctricos--- la ley de Ohm covariante es $j^\mu=\sigma E^\mu$, lo que
combinado con la descomposición da la cuadricorriente completa:
$J^\mu=\rho_e u^\mu+\sigma F^{\mu\nu}u_\nu$. La diferencia física fundamental:
en el límite $\sigma\to\infty$, para evitar que $j^\mu$ diverja y se destruya
la conservación de energía, $E^\mu$ debe anularse. Esto produce el
congelamiento topológico de $\mathbf{B}$. En RRMHD con $\sigma$ finita,
$E^\mu\neq 0$ permite tres cosas físicas reales: disipación resistiva con
calentamiento de Joule, difusión del campo magnético, y cambios topológicos
vía reconexión. Conviene recordar que la conductividad escalar $\sigma$ es
válida cuando la frecuencia de colisión domina sobre la ciclotrónica; en
campos muy intensos $\sigma$ se vuelve un tensor con corrientes de Hall.
\end{bryorador}

\begin{bryorador}{18}{La Ley de Ohm Relativista II}{1:20}{25:35}
Esta diapositiva concreta la proyección $3+1$ de la ley de Ohm en variables
eulerianas, observables directamente en el laboratorio. Usamos
$u^\mu=(\Gamma,\Gamma\mathbf{v})$. El campo eléctrico comóvil expresado en
variables del laboratorio es
$E^\mu=\bigl(\Gamma(\mathbf{E}\cdot\mathbf{v}),\,\Gamma(\mathbf{E}+
\mathbf{v}\times\mathbf{B})\bigr)$. Tomando la componente temporal de
$J^\mu$ obtenemos la relación entre la carga laboratorio $\rho_q$ y la carga
propia $\rho_e$, con una corrección $\sigma(\mathbf{E}\cdot\mathbf{v})$. La
componente espacial nos da
$\mathbf{J}=\rho_e\Gamma\mathbf{v}+\sigma\Gamma(\mathbf{E}+\mathbf{v}\times
\mathbf{B})$. Sustituyendo $\rho_e$ y simplificando llegamos a la forma
cerrada:
\[
\mathbf{J}=\rho_q\mathbf{v}+\sigma\Gamma\bigl[\mathbf{E}+\mathbf{v}\times
\mathbf{B}-(\mathbf{E}\cdot\mathbf{v})\mathbf{v}\bigr].
\]
La interpretación física es muy rica: el primer término es arrastre
convectivo; el segundo, conducción óhmica relativista; el tercero, término
dinamoeléctrico $\mathbf{v}\times\mathbf{B}$; el cuarto, acoplamiento
longitudinal puramente relativista; y el factor $\Gamma$ global captura la
contracción espacial sobre el transporte conductivo. Ésta es la corriente que
efectivamente alimenta la fuente del sistema conservativo.
\end{bryorador}

\begin{bryorador}{19}{El Sistema Conservativo Completo}{1:35}{27:10}
Cierro mi bloque teórico mostrando el sistema en su forma canónica
hiperbólica: $\partial_t \mathbf{U}+\partial_i \mathbf{F}^i(\mathbf{U})=
\mathbf{S}(\mathbf{U})$. Éste es el formato que reciben los solucionadores
tipo Godunov. Los flujos espaciales $\mathbf{F}^i$, en el panel central,
contienen los términos de transporte advectivo: $Dv^i$ en la primera fila,
$S^j v^i+p\delta^{ij}$ para el momento, $(\tau+p)v^i$ para la energía, y las
contribuciones cruzadas $\epsilon^{ijk}E_k$ y $-\epsilon^{ijk}B_k$ para los
campos. Las componentes $\Psi$ y $\Phi$ aparecen como términos diagonales
adicionales, garantizando que las restricciones de divergencia se propaguen
hiperbólicamente.

El vector fuente, a la derecha, recoge los términos no homogéneos: la fuerza
de Lorentz $\rho_q\mathbf{E}+\mathbf{J}\times\mathbf{B}$ en la ecuación de
momento, el calentamiento $\mathbf{J}\cdot\mathbf{E}$ en la ecuación de
energía, la corriente $-\mathbf{J}$ en la ecuación de Faraday, y los
términos de amortiguamiento $-\kappa_B\Psi$ y $\rho_q-\kappa_E\Phi$. Los
términos en color naranja son las correcciones de Godunov--Powell, que
compensan la aceleración artificial de monopolos espurios y garantizan
localmente la segunda ley de la termodinámica. La implementación numérica
combina solucionador HLLC para la advección, esquemas IMEX Runge--Kutta para
los términos rígidos, y limpieza GLM hiperbólica--parabólica para las
divergencias.
\relevo{Le paso la palabra a Sebastián para una diapositiva crítica sobre IMEX.}
\end{bryorador}

%==============================================================================
\section{IMEX-RK --- Diapositiva 20 (Sebastián)}
%==============================================================================

\begin{seborador}{20}{Por Qué IMEX-RK es Indispensable}{1:35}{28:45}
Tomo esta diapositiva porque es donde se entiende por qué la elección
numérica importa tanto en RRMHD. El sistema tiene una rigidez intrínseca: en
el límite de alta conductividad, el término fuente del campo eléctrico ---que
va como $-\sigma\Gamma\mathbf{E}$--- tiene una escala temporal característica
$\tau_\text{stiff}\sim 1/\sigma$. Para $\sigma=10^4$ y un paso advectivo
$\Delta t_\text{adv}\sim 10^{-3}$, el término rígido exigiría $\Delta
t_\text{expl}\lesssim 10^{-4}$, es decir, un orden de magnitud más caro que
la advección. Con un esquema explícito puro esto es prohibitivo: las
simulaciones se vuelven inviables.

La solución son los esquemas IMEX, implícitos--explícitos, propuestos para
RRMHD por Palenzuela y por Aloy y Cordero-Carrión. Tratan implícitamente sólo
los términos rígidos: $\mathbf{E}^{n+1}$ aparece a ambos lados de la
ecuación, lo que permite resolver una ecuación lineal en cada celda. Los
términos no rígidos, advección y limpieza GLM, se siguen tratando
explícitamente. El esquema queda en tres etapas: primero avanzo los términos
explícitos con Runge--Kutta clásico; segundo, resuelvo implícitamente el
bloque rígido del campo eléctrico; tercero, actualizo las variables
conservativas a primitivas. El resultado es estabilidad para cualquier
$\sigma$, incluyendo el límite cuasi-ideal $\sigma\to\infty$. Sin IMEX la
simulación se paralizaría por agotamiento del paso de tiempo.
\relevo{Le paso la palabra a Bryan para entrar en la inestabilidad.}
\end{seborador}

%==============================================================================
\section{KHI: descripción física --- Diapositiva 21 (Bryan)}
%==============================================================================

\begin{bryorador}{21}{Descripción Física General de la KHI}{1:20}{30:05}
Gracias, Sebastián. Entramos al capítulo tres, donde describimos la
inestabilidad Kelvin--Helmholtz desde lo más físico hasta el caso RRMHD. El
mecanismo básico, que se remonta a los teoremas originales de Kelvin y
Helmholtz, ocurre en la interfaz entre dos capas de fluido o plasma con
velocidades relativas distintas. Cinéticamente: la diferencia de velocidades
induce variaciones locales de presión que empujan la interfaz alejándola del
equilibrio, lo que produce el crecimiento exponencial característico.

La evolución morfológica tiene cuatro fases: fase lineal, donde las
perturbaciones crecen como $e^{\gamma t}$; saturación, con deformación severa
de la interfaz; enrollamiento, donde la capa de vorticidad se enrosca
formando los vórtices de Kelvin--Helmholtz; y cascada, con estructuras
secundarias y eventual turbulencia. Los vórtices KH convierten energía
cinética del flujo de cizallamiento en energía magnética, amplificando
$\mathbf{B}$. La figura muestra el perfil de velocidad
$V_y(x)=v_\text{sh}\tanh(x/a_\text{kh})$ con $v_\text{sh}=0.5$ y
$a_\text{kh}=0.05$, que es el setup de la campaña. La relevancia astrofísica
es triple: transporte cruzado de momento y energía, mezcla física entre
capas, reconexión turbulenta; y por su simplicidad geométrica, problema test
estándar para validar códigos RRMHD.
\relevo{Le devuelvo la palabra a Sebastián para la deducción clásica.}
\end{bryorador}

%==============================================================================
\section{KHI: clásica, MHD, RMHD --- Diapositivas 22--24 (Sebastián)}
%==============================================================================

\begin{seborador}{22}{KHI en Hidrodinámica Clásica: Análisis Lineal}{1:35}{31:40}
Tomo el caso clásico, hidrodinámica incompresible y sin gravedad.
Consideramos dos fluidos con densidades $\rho_1$ y $\rho_2$, separados por
una interfaz horizontal y moviéndose con velocidades horizontales $U_1$ y
$U_2$. Como ambos son incompresibles, los potenciales de velocidad satisfacen
Laplace. Ensayamos una perturbación de la interfaz
$\eta=\hat\eta\,e^{i(kx-\omega t)}$, y aplicamos dos condiciones de contorno:
la cinética, que dice que la interfaz se mueve con el fluido; y la dinámica,
que exige continuidad de la presión a través de la interfaz. Estas dos
condiciones combinadas producen una ecuación cuadrática en $\omega$.

Lo que hace especial a la KHI clásica es que el discriminante de esta
cuadrática es siempre negativo cuando $U_1\neq U_2$: explícitamente,
$\Delta=-4k^2\rho_1\rho_2(U_1-U_2)^2$. Esto fuerza a las raíces de $\omega$ a
ser complejas conjugadas, con parte imaginaria positiva. Como $e^{-i\omega t}
\propto e^{\operatorname{Im}(\omega)t}$, la perturbación crece
exponencialmente para cualquier longitud de onda. La conclusión es
contundente: en hidrodinámica clásica la cizalladura genera inestabilidad
\emph{incondicional}. Cualquier modo, por largo o corto que sea, es inestable
mientras haya diferencia de velocidad y no exista gravedad ni tensión
superficial. Ésta es la versión más pura del fenómeno.
\end{seborador}

\begin{seborador}{23}{KHI en MHD: el Campo como Fuerza Restauradora}{1:35}{33:15}
Cuando agregamos un campo magnético uniforme paralelo al flujo, $\mathbf{B}_0
=B_0\hat{i}$, la situación cambia cualitativamente. De la ecuación de
inducción linealizada en el límite ideal, donde se cumple congelamiento,
despejamos la perturbación del campo magnético en función de la perturbación
de velocidad. Sustituyendo en la presión total perturbada, magnética más
cinética, obtenemos un término extra de tensión magnética. Aplicando
continuidad de presión total en la interfaz, la relación de dispersión MHD se
reduce a $\rho_1(\omega-kU_1)^2+\rho_2(\omega-kU_2)^2=2k^2 B_0^2/\mu_0$. El
término derecho es la \emph{tensión magnética} resistiendo la deformación.

La condición de inestabilidad ahora exige que la energía cinética de cizalla
supere la tensión magnética: el discriminante negativo requiere
$\rho_1\rho_2(U_1-U_2)^2>2B_0^2(\rho_1+\rho_2)/\mu_0$. La inestabilidad ya no
es incondicional. Un campo paralelo suficientemente fuerte la suprime; un
campo perpendicular permite desarrollo más libre, como discutió Bucciantini
y Del Zanna. Hacia el RRMHD: cuando rompemos el congelamiento por $\sigma$
finita, modos que la MHD ideal suprimiría pueden volver a crecer por
difusión del campo, abriendo nuevas regiones de inestabilidad.
\end{seborador}

\begin{seborador}{24}{KHI en RMHD Relativista}{1:25}{34:40}
Cerrando el bloque, paso al caso magnetohidrodinámico relativista. La
relatividad importa cuando las velocidades del flujo se aproximan a $c$, o
cuando la energía interna específica es comparable a la energía de masa en
reposo. En RMHD la inercia efectiva del fluido aumenta por contribuciones de
presión térmica y entalpía: $\rho h\Gamma^2>\rho\Gamma^2$. Coloquialmente
decimos que el fluido «pesa más». El efecto sobre la tasa de crecimiento fue
estudiado por Osmanov y colaboradores en 2008, mostrando que la inercia
efectiva incrementada tiende a reducir $\gamma_\text{KHI}$ comparado con el
caso newtoniano, estabilizando los modos de onda larga. Pimentel y
Lora-Clavijo en 2019 encontraron además que la orientación del campo
$\mathbf{B}$ y el factor de Lorentz $\Gamma$ modifican significativamente los
dominios de inestabilidad.

La condición de inestabilidad relativista se expresa en términos del cociente
entre la velocidad de Alfvén relativista $V_A=B_0/\sqrt{\mu_0\rho h}$ y la
velocidad del sonido. El umbral disminuye al aumentar $\Gamma$. Esto tiene
aplicaciones directas: en nebulosas de viento de púlsar la KHI modula la
emisión sincrotrón, y en chorros de núcleos galácticos activos la KHI
determina si el chorro permanece colimado o se fragmenta.
\relevo{Le paso la palabra a Bryan para el caso RRMHD y los diagnósticos.}
\end{seborador}

%==============================================================================
\section{KHI en RRMHD y diagnósticos --- Diapositivas 25--27 (Bryan)}
%==============================================================================

\begin{bryorador}{25}{KHI en RRMHD: Impacto Dual de la Resistividad}{1:35}{36:15}
Gracias, Sebastián. Tocamos el caso RRMHD, que es el de nuestra tesis. La
resistividad rompe el congelamiento: la ley de Ohm con $\eta=1/\sigma$
introduce un campo eléctrico comóvil distinto de cero, lo que permite
\emph{difusión} del campo magnético a través del fluido. Éste es el mecanismo
físico que abre el escenario nuevo.

La resistividad tiene un efecto \emph{dual} sobre la KHI. Como amortiguador,
actúa como mecanismo disipativo para los modos de alta frecuencia, reduciendo
$\gamma_\text{KHI}$ en el régimen resistivo fuerte: el campo difunde antes de
que la inestabilidad pueda crecer. Como activador, en cambio, elimina
restricciones topológicas de la MHD ideal, permitiendo inestabilidades tipo
\emph{tearing} que pueden coexistir o acoplarse a la KHI; en ciertos
regímenes la reconexión puede incluso acelerar el crecimiento. Hay también
una rigidez numérica inducida que justifica la elección IMEX que ya discutió
Sebastián.

El resultado de nuestra campaña confirma la hipótesis estabilizadora
dominante: cuando aumentamos $\sigma$, $\gamma_\text{KHI}$ tiende a su valor
asintótico $\gamma_0=0.894$; cuando bajamos $\sigma$, $\gamma_\text{KHI}$
tiende a cero. El ajuste sigmoide v5 cuantifica el centro de la transición en
$\sigma_0=4472\pm 81$. Éste es el primer estimador paramétrico exacto de
$\sigma_\text{crit}$ que producimos.
\end{bryorador}

\begin{bryorador}{26}{Vorticidad y Enstrofia: Diagnósticos}{1:30}{37:45}
Esta diapositiva define los diagnósticos cuantitativos del desarrollo de la
KHI. El perfil de velocidad de cizalladura es
$V_y(x)=v_\text{sh}\tanh\bigl((x-x_0)/a_\text{kh}\bigr)$. La vorticidad
relevante en dos dimensiones es la componente $z$:
$\omega_z=\partial_x V_y-\partial_y V_x$. La enstrofía total se define como
la integral del cuadrado de la vorticidad. Cuando incluimos la posibilidad de
movimiento fuera del plano, calculamos la fracción $f_{Vz}$ que mide cuánto
se desvía el flujo del régimen estrictamente bidimensional; valores próximos
a cero indican que la dinámica permanece en el plano, condición que
verificamos en nuestras simulaciones.

El observable principal de la campaña es la \emph{enstrofía de la
perturbación}. La enstrofía total está dominada por la fuerte cizalladura del
flujo base; para aislar las estructuras inestables emergentes restamos a la
vorticidad instantánea su valor inicial promediado en $y$, obteniendo
$\omega'_z$. La integral del cuadrado de esta perturbación, $\Omega'_z(t)$,
es libre de la contaminación del fondo y permite identificar limpiamente las
fases de crecimiento lineal y saturación no lineal. Ésta es la cantidad cuyo
logaritmo crece linealmente en la fase exponencial y la que ajustamos para
extraer $\gamma_\text{KHI}$. Sin esta sustracción, la señal de la
inestabilidad quedaría enterrada en la cizalladura de fondo.
\end{bryorador}

\begin{bryorador}{27}{Cálculo de la Tasa de Crecimiento Lineal}{1:30}{39:15}
En esta diapositiva detallo la metodología numérica robusta para extraer
$\gamma_\text{KHI}$. La base teórica es la ya mencionada: durante la fase
lineal la amplitud crece exponencialmente y, dado que $\Omega'_z\propto A^2$,
$\ln\Omega'_z(t)=C+2\gamma_\text{KHI}t$. Hacemos un ajuste lineal en una
ventana de crecimiento limpio, con pendiente $s=2\gamma_\text{KHI}$, y por
tanto $\gamma_\text{KHI}=s/2$.

Los pasos del análisis son cinco: primero, graficamos $\Omega'_z(t)$ e
identificamos visualmente la fase lineal; segundo, graficamos el logaritmo y
seleccionamos una ventana fija $t\in[2.4,3.4]$, idéntica para todas las
simulaciones para evitar sesgo; tercero, ajustamos el modelo lineal por
mínimos cuadrados; cuarto, calculamos $\gamma_\text{KHI}$; y quinto,
repetimos el procedimiento para los cuarenta y cuatro valores de $\sigma$. La
gráfica a la derecha muestra $\ln\Omega'_z(t)$ para todas las simulaciones,
con la ventana lineal destacada en amarillo. La bondad de ajuste promedio es
$\bar R^2=0.92$, y el $R^2_\text{mín}=0.31$ se usa como filtro de
supervivencia. El procedimiento validó las predicciones de Pimentel--Lora y
Osmanov, y aportó adicionalmente un estimador paramétrico exacto de
$\sigma_\text{crit}$ vía el modelo logístico.
\relevo{Le paso la palabra a Sebastián para el balance editorial.}
\end{bryorador}

%==============================================================================
\section{Avance editorial --- Diapositivas 28--29 (Sebastián)}
%==============================================================================

\begin{seborador}{28}{Estructura y Avance de Escritura I}{1:30}{40:45}
Gracias, Bryan. Tomo la palabra para informar el estado de la monografía. El
capítulo uno, introducción, está completo y revisado, con aproximadamente
catorce páginas. Cubre el contexto de plasmas relativistas, la justificación
de RRMHD frente a RMHD ideal, y los objetivos del trabajo. El capítulo dos,
marco teórico, también está completo, con cuarenta y cinco páginas. Su
estructura corresponde directamente a las diapositivas que acabamos de
presentar: electrodinámica covariante con tensor de Faraday, leyes de
Maxwell, derivación variacional del tensor energía--momento, sistema GLM
aumentado, ley de Ohm relativista en proyección $3+1$, sistema conservativo
de trece componentes, análisis lineal de la KHI, y vorticidad con enstrofía
y cálculo de $\gamma_\text{KHI}$. El capítulo tres, KHI en los tres marcos,
está completo en veintidós páginas, cubriendo la descripción física general,
la derivación clásica, la extensión MHD, las correcciones relativistas RMHD,
y el rol de la resistividad en RRMHD con la escala efectiva $\ell_\text{eff}$.

Adicionalmente, hay actualizaciones activas en el marco teórico que se están
incorporando al capítulo dos: el modelo sub-escala v4 con $\gamma_0$ libre,
el modelo sigmoide v5 con saturación logística, los cuatro estimadores
independientes de $\sigma_\text{crit}$, y la discrepancia entre el criterio
lineal y el criterio global. Estas adiciones suman aproximadamente cinco
páginas más al capítulo dos.
\end{seborador}

\begin{seborador}{29}{Estructura y Avance de Escritura II}{1:35}{42:20}
Continuando con el balance, los capítulos en redacción activa son tres. El
capítulo cuatro, diseño numérico y setup, en progreso, cubre las condiciones
iniciales, la malla, los parámetros, la validación, las cuarenta y cuatro
simulaciones efectivas y la ventana lineal $t\in[2.4,3.4]$. El capítulo
cinco, resultados, también en progreso, contiene $\gamma_\text{KHI}(\sigma)$,
la enstrofía $\Omega'_z(t)$, la ley de potencias del pico no lineal, los
tiempos de saturación, el modelo sigmoide v5, $\sigma_\text{crit}=4472\pm 81$,
y la escala efectiva $\ell_\text{eff}\approx 0.27\,a_\text{kh}$. El capítulo
seis, discusión y conclusiones, queda como pendiente para la fase final,
donde compararemos sistemáticamente los criterios lineal y global,
discutiremos la discrepancia teoría--simulación, las limitaciones del trabajo
y las perspectivas. La tabla resume el estado: el total actual es de
aproximadamente noventa y cinco páginas de contenido.

La actualización inmediata del marco teórico es la inclusión, en el capítulo
dos, de la ecuación de evolución de la vorticidad en RRMHD, mostrada abajo a
la derecha. Esta ecuación, derivada del rotacional de la ecuación de momento,
tiene un término inercial $\nabla\times(\mathbf{v}\times\boldsymbol\omega)$ y
un término magnético dependiente del rotacional de
$\mathbf{J}\times\mathbf{B}/(\rho h)$. Justifica formalmente por qué
$\Omega'_z$ es el observable adecuado para diagnosticar la KHI.
\relevo{Le paso la palabra a Bryan para la campaña de simulaciones.}
\end{seborador}

%==============================================================================
\section{Campaña numérica --- Diapositivas 30--32 (Bryan)}
%==============================================================================

\begin{bryorador}{30}{Diseño de la Campaña Numérica I: Setup y Parámetros}{1:25}{43:45}
Gracias, Sebastián. Entramos al capítulo de la campaña numérica. El setup del
código \textsc{cueva} usa el perfil de cizalladura tangente hiperbólico
estándar con $v_\text{sh}=0.5$ y $a_\text{kh}=0.05$. El campo magnético tiene
polaridades opuestas a cada lado de la interfaz, formando una lámina de
corriente. El dominio es bidimensional periódico, con resolución variable
para estudios de convergencia. El observable principal es la enstrofía de la
perturbación que ya hemos discutido.

Las estadísticas globales de la campaña son: cuarenta y cuatro simulaciones
aprobadas tras filtros de calidad; rango de conductividad cubriendo desde
$10^2$ hasta $1.05\times 10^4$, es decir, tres órdenes de magnitud completos;
bondad de ajuste promedio $\bar R^2=0.92$; y un mínimo de $R^2_\text{mín}=0.31$
como filtro de supervivencia, que descartó simulaciones donde el ajuste lineal
era ruidoso. La distribución del barrido en $\sigma$ fue intencionalmente
densificada en la zona de transición, entre $2.8\times 10^3$ y
$1.05\times 10^4$, donde se identifica la transición resistivo--cuasi-ideal,
en lugar del barrido grueso original de siete valores del póster. La tabla
muestra un extracto representativo de diez puntos: para cada $\sigma$, el
pico de enstrofía $\Omega^\text{máx}_{zp}$, el tiempo $t_\text{peak}$, y el
número de Reynolds magnético modificado $Rm_*$.
\end{bryorador}

\begin{bryorador}{31}{Visualización del Setup: Campo Vectorial}{1:15}{45:00}
Esta diapositiva es la materialización gráfica del setup que acabo de
describir. Lo que ven en pantalla es el campo de velocidad bidimensional,
$|\vec V(x,y,t=0)|$, en el tiempo cero de la simulación. La paleta de color
es \emph{inferno}: las regiones más brillantes corresponden a mayor velocidad,
las más oscuras a menor velocidad. Las líneas blancas son \emph{streamlines},
las trayectorias instantáneas del fluido.

Lo importante de esta figura es lo siguiente: aparecen \emph{dos} láminas de
cizalladura simétricas, una en $x=-0.5$ y otra en $x=+0.5$. Esto no es un
artefacto numérico, es una consecuencia directa de la condición de frontera
periódica del dominio: para que el campo de velocidad sea continuo en los
bordes, debe existir una segunda cizalladura que cierre el patrón de manera
consistente. El perfil tangente hiperbólico que aplicamos analíticamente sólo
introduce una cizalladura, pero la periodicidad del dominio induce su imagen
especular. En cada una de estas dos láminas se desarrolla, en paralelo, una
inestabilidad Kelvin--Helmholtz independiente. Esto duplica nuestra
estadística por simulación y permite verificar la simetría del esquema
numérico: si las dos cizalladuras evolucionaran asimétricamente, sabríamos
que hay un problema. En todas nuestras simulaciones la simetría se conserva
con altísima fidelidad.
\end{bryorador}

\begin{bryorador}{32}{Diseño de la Campaña Num.\ II: Modelo Sigmoide v5}{1:50}{46:50}
En esta diapositiva consolidamos los diagnósticos cuantitativos y el modelo
sigmoide. Implementamos cuatro diagnósticos independientes de
$\sigma_\text{crit}$. Primero, $\gamma_\text{KHI}(\sigma)$ vía el ajuste
lineal del logaritmo de la enstrofía. Segundo, el pico no lineal
$\Omega^\text{máx}_{zp}(\sigma)$. Tercero, $t_\text{peak}(\sigma)$. Y cuarto,
$\Omega^\text{máx}_{zp}(Rm_*)$ con $Rm_*=v_\text{sh}\,a_\text{kh}\,\sigma$.

El modelo sigmoide v5 propone $\gamma_\text{KHI}(\sigma)=\gamma_0/
\bigl(1+\exp[-k(\sigma-\sigma_0)]\bigr)$. Resuelve la patología de
no-identificabilidad del modelo sub-escala v4 previo, donde $\gamma_0$
aparecía en una posición que requería observar la saturación, lo cual no
era posible en todas las simulaciones. En el sigmoide v5, $\gamma_0$ aparece
como factor multiplicativo y es identificable sin requerir saturación. Los
parámetros ajustados con $N=44$ son: $\gamma_0=0.8941\pm 0.0119$;
$\sigma_0=4471.94\pm 80.56$; $k=(4.97\pm 0.13)\times 10^{-4}$; y $R^2=0.9963$.

La gráfica muestra los datos con los ajustes superpuestos y los cuatro
estimadores marcados con líneas verticales. La tabla comparativa abajo
indica que pasamos de $R^2=0.9733$ con el modelo clásico a $R^2=0.9963$ con
el sigmoide, ganando algo más de dos puntos porcentuales que en este
contexto son significativos. El ajuste sigmoide identifica
$\sigma_\text{crit}^\text{sigm}\equiv\sigma_0$ como el primer estimador
paramétrico \emph{exacto} de la conductividad crítica.
\relevo{Le paso la palabra a Sebastián para los regímenes y resultados.}
\end{bryorador}

%==============================================================================
\section{Tres regímenes y resultados --- Diapositivas 33--35 (Sebastián)}
%==============================================================================

\begin{seborador}{33}{Los Tres Regímenes Identificables por $\sigma_\text{crit}$}{1:50}{48:40}
Gracias, Bryan. Tomo la palabra para discutir la interpretación física. La
campaña permite identificar tres regímenes claramente diferenciados por el
valor de $\sigma$ respecto a $\sigma_\text{crit}$.

En el régimen \emph{resistivo fuerte}, $\sigma\ll\sigma_\text{crit}$, la
difusión magnética domina sobre la advección. La inestabilidad queda
fuertemente amortiguada porque el campo difunde antes de que pueda crecer
macroscópicamente: $\gamma_\text{KHI}\to 0$, el pico de enstrofía es bajo,
$t_\text{peak}$ tardío, y no se observan vórtices coherentes. En nuestra
campaña este régimen se identifica para $\sigma\lesssim 10^3$.

En el régimen de \emph{transición}, $\sigma\sim\sigma_\text{crit}$, hay
competencia entre advección y difusión: aparecen cambios rápidos de pendiente
en $\Omega^\text{máx}_{zp}(\sigma)$ y en $t_\text{peak}(\sigma)$. En esta
zona pueden coexistir inestabilidades tipo \emph{tearing} acopladas a la KHI,
con reconexión incipiente. Es la zona más interesante físicamente, y es
justamente donde densificamos el barrido. El centro de la transición es
$\sigma_\text{crit}=4472\pm 81$.

En el régimen \emph{cuasi-ideal}, $\sigma\gg\sigma_\text{crit}$, la advección
domina y la KHI crece libremente, aproximándose al régimen RMHD ideal.
$\gamma_\text{KHI}\to\gamma_0=0.894$, el pico es alto, $t_\text{peak}$
temprano, y los vórtices de Kelvin--Helmholtz se forman claramente.
Identificamos este régimen para $\sigma\gtrsim 7\times 10^3$. La predicción
operativa es robusta porque está sustentada por cuatro estimadores
independientes que coinciden en orden de magnitud.
\end{seborador}

\begin{seborador}{34}{Resultados Cuantitativos: Ley de Potencias y Escala Efectiva}{2:00}{50:40}
Esta diapositiva consolida los hallazgos más fuertes de la tesis. La primera
ley de potencias que ajustamos es para el pico no lineal:
$\Omega^\text{máx}_{zp}\propto\sigma^\alpha$. El ajuste completo, con $N=29$
puntos limpios, da $\alpha=1.213$ con $R^2=0.997$. El ajuste asintótico
restringido a $\sigma\geq 6000$ da $\alpha=1.238$ con $R^2=0.993$. Esto es
físicamente importante porque la predicción teórica difusiva ingenua daría
$\alpha\approx 1/2$. El valor observado, alrededor de $1.2$, indica que el
control asintótico no es estrictamente difusivo: involucra estructuras
secundarias no lineales ---láminas de corriente y vórtices secundarios---
que aceleran la cascada.

La escala efectiva $\ell_\text{eff}=(\gamma_0\sigma)^{-1/2}\approx 0.0137$,
equivalente a $0.27\,a_\text{kh}$, es clave: la transición está controlada
por escalas internas \emph{menores} que el espesor nominal de la capa de
corte. Esto explica naturalmente la discrepancia entre los criterios. La
tabla contrasta los cuatro estimadores experimentales ---6800, 3400, 6000 y
$4472\pm 81$--- con los tres estimadores teóricos: 447 del criterio lineal,
490 con corrección magnética, y 8595 del criterio de fracción $f=0.6$. La
discrepancia entre el $\sigma_\text{crit}$ global experimental y el
$\sigma_\text{crit}$ lineal teórico es de un orden de magnitud, factor de 10
a 15. Esto, lejos de ser una inconsistencia, es la principal contribución
conceptual del trabajo: el criterio lineal mide el umbral local a escala
$a_\text{kh}$, mientras que el criterio global mide el encendido no lineal a
la escala efectiva $\ell_\text{eff}\ll a_\text{kh}$, asociada a la formación
de hojas de corriente y vórtices secundarios.
\end{seborador}

\begin{seborador}{35}{Perfiles de Vorticidad: Espacio de Parámetros}{1:25}{52:05}
Cierro mi último bloque mostrando los perfiles de vorticidad. El campo de
vorticidad inicial está determinado exclusivamente por el perfil de
cizalladura: $\omega^{(0)}_z(x)=\partial_x V_y=(v_\text{sh}/a_\text{kh})\,
\operatorname{sech}^2(x/a_\text{kh})$. Esto es un pico estrecho centrado en
$x=0$, con ancho aproximado $2a_\text{kh}=0.1$, y amplitud
$v_\text{sh}/a_\text{kh}=10$.

La evolución esperada de $\omega_z(x,y,t)$ pasa por cuatro fases: en $t=0$ es
el pico estrecho descrito; en el tiempo lineal aparece la modulación
periódica en $y$ correspondiente a los modos KH; en el tiempo de saturación
se forman estructuras vorticales coherentes; y en el régimen no lineal
aparece el enrollamiento y la cascada turbulenta. La resistividad difumina y
debilita el pico de vorticidad, lo que se traduce directamente en el
amortiguamiento que observamos en el régimen resistivo. En el espacio
reservado anuncio las figuras adicionales que vamos a incluir al consolidar
el capítulo cinco: campos de color $\omega_z(x,y,t)$ para varios tiempos y
distintos $\sigma$ con paleta inferno; perfil unidimensional
$\omega_z(0,y,t)$ contra $y$; líneas de campo magnético superpuestas; y
animaciones de la evolución temporal completa.
\relevo{Le paso la palabra a Bryan para el cierre.}
\end{seborador}

%==============================================================================
\section{Cierre --- Diapositivas 36--37 (Bryan)}
%==============================================================================

\begin{bryorador}{36}{Estado Actual del Proyecto}{1:25}{53:30}
Gracias, Sebastián. Cierro la presentación con el balance del proyecto a mayo
de 2026. En lo \emph{completado}: los capítulos uno, dos y tres están
redactados con un total aproximado de ochenta y un páginas; tenemos las
cuarenta y cuatro simulaciones aprobadas; el modelo sigmoide v5 con
$R^2=0.9963$; la ley de potencias $\Omega^\text{máx}_{zp}\propto\sigma^{1.213}$;
los cuatro estimadores independientes de $\sigma_\text{crit}$; la escala
efectiva $\ell_\text{eff}\approx 0.27\,a_\text{kh}$; y la predicción teórica
$\sigma^\text{lin}_\text{crit}=1/(\gamma_0 a_\text{kh}^2)$, contrastada con
los resultados numéricos.

En lo que está actualmente \emph{en progreso}: el capítulo cuatro de diseño
numérico detallado; el capítulo cinco de resultados cuantitativos
consolidados; la inclusión en el capítulo dos de la ecuación de evolución de
la vorticidad en RRMHD; la documentación formal de los modelos sigmoide v5 y
sub-escala v4; y los diagnósticos avanzados sugeridos: espectros de potencia
y calentamiento óhmico.

Como \emph{próximos pasos}: capítulo seis de discusión y conclusiones; las
simulaciones del cuarto objetivo específico, sobre orientación variable del
campo magnético; un par de simulaciones de alta resolución para acceder
limpiamente al régimen no lineal; y la defensa final, prevista para junio
de 2026.
\end{bryorador}

\begin{bryorador}{37}{Síntesis: Contribuciones y Perspectivas}{1:25}{54:55}
Termino con la síntesis de contribuciones y perspectivas. Las contribuciones
consolidadas del trabajo son ocho. Primero, una fundamentación rigurosa del
marco RRMHD desde formas diferenciales y acción variacional. Segundo, la
derivación explícita del sistema GLM con su origen covariante. Tercero, la
ley de Ohm relativista deducida con proyección $3+1$ explícita. Cuarto, el
sistema conservativo de trece componentes con las correcciones de
Godunov--Powell. Quinto, el análisis sistemático de la KHI desde lo clásico
hasta RRMHD, mostrando explícitamente cómo cambia la condición de
inestabilidad en cada nivel. Sexto, el modelo sigmoide v5 como primer
estimador paramétrico exacto de $\sigma_\text{crit}$. Séptimo, la campaña de
cuarenta y cuatro simulaciones cuantificadas y validadas con diagnósticos
múltiples. Y octavo, la identificación de la escala efectiva
$\ell_\text{eff}$ como controlador del encendido no lineal y su rol como
explicación de la discrepancia entre criterios lineal y global.

Los resultados cuantitativos clave son: $\sigma_\text{crit}=4472\pm 81$,
$\gamma_0=0.894\pm 0.012$, $\Omega^\text{máx}_{zp}\propto\sigma^{1.213}$, y
$\ell_\text{eff}\approx 0.27\,a_\text{kh}$. Como próxima fase, durante los
meses restantes del semestre, ejecutaremos: simulaciones con orientación
variable del campo magnético; espectros de potencia $E_k(k,t,\sigma)$; tasa
de calentamiento óhmico $\int J^2/\sigma\,dA$; cuantificación de la
magnetización $\sigma_m=b^2/(\rho h)$; y comparación sistemática con las
configuraciones de Mizuno (2013) y Pimentel--Miranda (2019). Con esto
cerramos. \emph{Muchas gracias por su atención.} \escena{pausa, gesto de
agradecimiento; abrir turno de preguntas}
\end{bryorador}

%==============================================================================
\vspace{1em}
\begin{center}
\begin{tcolorbox}[colback=tagbg, colframe=neutral!50, width=0.9\linewidth,
                  arc=2pt, boxrule=0.5pt]
\small\centering
\textbf{Tiempo total estimado:} $\sim 55$ minutos a ritmo conservador
($\sim 150$ wpm) \\
$\sim 45$ minutos a ritmo de defensa típico ($\sim 180$ wpm). \\[0.4em]
Ajustar densidad oral según ensayo previo. Recomendado cronometrar las
diapositivas 10, 13, 32 y 34, que son las más densas.
\end{tcolorbox}
\end{center}

\end{document}